\begin{figure}[H]
\centering
\begin{tikzpicture}[font=\small,>=Latex]
  \begin{scope}[scale=1.12]
    \fill[blue!5] (-1.7,-1.7) rectangle (1.7,1.7);
    \fill[orange!18]
      (0,1.35) .. controls (.80,1.35) and (1.02,.95) .. (1.02,.60)
      -- (1.02,-.60)
      .. controls (1.02,-.95) and (.80,-1.35) .. (0,-1.35)
      .. controls (-.80,-1.35) and (-1.02,-.95) .. (-1.02,-.60)
      -- (-1.02,.60)
      .. controls (-1.02,.95) and (-.80,1.35) .. (0,1.35);
    \draw[thick]
      (0,1.35) .. controls (.80,1.35) and (1.02,.95) .. (1.02,.60)
      -- (1.02,-.60)
      .. controls (1.02,-.95) and (.80,-1.35) .. (0,-1.35)
      .. controls (-.80,-1.35) and (-1.02,-.95) .. (-1.02,-.60)
      -- (-1.02,.60)
      .. controls (-1.02,.95) and (-.80,1.35) .. (0,1.35);
    \draw[gray] (-1.7,-1.7) rectangle (1.7,1.7);
    \draw[->,gray] (-1.87,0)--(1.91,0) node[right,black] {\(x\)};
    \draw[->,gray] (0,-1.87)--(0,1.91) node[above,black] {\(y\)};
    \fill[white] (-.25,-.25) rectangle (.25,.25);
    \draw[thick] (-.25,-.25) rectangle (.25,.25);
    \node at (.02,.01) {\(\mathcal P\)};
    \node at (.0,.76) {\(K<0\)};
    \node at (1.28,1.42) {\(K>0\)};
    \node[fill=white,inner sep=1pt] at (.90,-1.15) {\(K=0\)};
  \end{scope}
\end{tikzpicture}
\caption{The model patch, schematically and not to scale.
The curvature zero set lies outside the perturbation box
\(\mathcal P\). The positive exterior permits estimates on lower and
upper caps whose free artificial edges may cut through the negative
region. The lower-cap integration region is shown in \Cref{fig:cap}.}
\label{fig:geometry}
\end{figure}
